% Titolo: Revisione documenti per candidatura ed elezione responsabile
\documentclass[a4paper, 11pt]{article}

% --- PACCHETTI ---
\usepackage[utf8]{inputenc}
\usepackage[italian]{babel}
\usepackage[margin=2.5cm]{geometry}
\usepackage{graphicx}
\usepackage{tabularx}
\usepackage{hyperref}
\usepackage{fancyhdr}
\usepackage{array}
\usepackage{titlesec}
\usepackage{xcolor}

% --- CONFIGURAZIONE HEADER E FOOTER ---
\pagestyle{fancy}
\fancyhf{}
\fancyhead[L]{Bit Crusher}
\fancyhead[R]{Verbale Interno N. 5} %INSERIRE NUMERO RIUNIONE 
\fancyfoot[C]{\thepage}
\renewcommand{\headrulewidth}{0.4pt}

% Configurazione link indice
\hypersetup{
    colorlinks=true,
    linkcolor=black,
    urlcolor=blue
}

\begin{document}

% ==========================================
% FRONTESPIZIO
% ==========================================
\begin{titlepage}
    \centering
    
    % --- LOGHI ---
    % Sostituisci i \framebox con i comandi \includegraphics corretti
    \begin{minipage}{0.45\textwidth}
        \centering
        % Esempio: \includegraphics[width=0.7\textwidth]{logo_unipd.png}
        \includegraphics[width=0.7\textwidth]{unipd_logo.png}
    \end{minipage}
    \hfill
    \begin{minipage}{0.45\textwidth}
        \centering
        % Esempio: \includegraphics[width=0.7\textwidth]{WhatsApp Image 2026-10-15 at 17.29.12.jpeg}
        \includegraphics[width=0.7\textwidth]{Logo-BitCrusher-Nero.png}
    \end{minipage}

    \vspace{1.5cm}
    {\Large \textbf{UNIVERSITÀ DEGLI STUDI DI PADOVA}} \\
    \vspace{0.3cm}
    {\large Corso di Laurea in Informatica} \\
    \vspace{0.3cm}
    {\large Insegnamento di Ingegneria del Software} \\
    \vspace{0.3cm}
    {\large Anno Accademico 2026/2027} \\

    \vspace{2cm}
    {\Huge \textbf{Verbale di Riunione}} \\
    \vspace{0.5cm}
    {\Large Riunione Interna \#5}  %INSERIRE NUMERO RIUNIONE

    \vspace{2cm}
    \begin{flushleft}
        \large
        \textbf{Gruppo:} Bit Crusher \\
        \textbf{Email:} bitcrusherteam@gmail.com \\
        \textbf{Versione:} v1.0.0 \\
        \textbf{Data della riunione:} 2026-10-29 \\ 
        \textbf{Data ultima modifica:} 29 ottobre 2026
    \end{flushleft}
    \vfill
\end{titlepage}

\newpage

% ==========================================
% REGISTRO DELLE MODIFICHE
% ==========================================
\section*{Registro delle Modifiche}
\vspace{0.5cm}
\begin{center}
\renewcommand{\arraystretch}{1.5}
\begin{tabularx}{\textwidth}{|l|l|l|l|X|}
\hline
\textbf{Versione} & \textbf{Data} & \textbf{Autore} & \textbf{Verificatore} & \textbf{Descrizione} \\ \hline
1.0.0 & 2026-11-03 & Alex Singh & - & Approvazione per il rilascio. \\ \hline
0.2.0 & 2026-11-03 & Antonio Tufano & Riccardo Pozzato & Aggiornamento logo UniPD e adeguamento registro modifiche. \\ \hline
0.1.0 & 2026-10-29 & Alexandru Mitu & Pietro Leonardo Acampora & Prima stesura verbale \\ \hline
\end{tabularx}
\end{center}

\newpage

% ==========================================
% INDICE
% ==========================================
\tableofcontents
\newpage

% ==========================================
% 1. INFORMAZIONI GENERALI
% ==========================================
\section{Informazioni Generali}
\begin{itemize}
    \item \textbf{Luogo/Piattaforma:} Riunione Discord
    \item \textbf{Data:} 2026-10-29
    \item \textbf{Orario di inizio:} 17:40
    \item \textbf{Orario di fine:} 20:30
    \item \textbf{Responsabile:} Alex Singh
    \item \textbf{Scriba:} Pietro Acampora
\end{itemize}

\vspace{0.5cm}
\textbf{Partecipanti presenti:} \\
\vspace{0.2cm}
\renewcommand{\arraystretch}{1.2}
\begin{tabularx}{0.8\textwidth}{|X|c|}
\hline
\textbf{Nome e Cognome} & \textbf{Matricola} \\ \hline
Pietro Leonardo Acampora & 2146762 \\ \hline
Marco Bollettin & 2147955 \\ \hline
Alexandru Mitu & 2101083 \\ \hline
Riccardo Pozzato & 2147988 \\ \hline
Alex Singh & 2148621 \\ \hline
Antonio Tufano & 2137972 \\ \hline
\end{tabularx}

\vspace{0.5cm}
\textbf{Partecipanti assenti:} \\
\vspace{0.2cm}
\renewcommand{\arraystretch}{1.2}
\begin{tabularx}{0.8\textwidth}{|X|c|}
\hline
\textbf{Erik Giuliano} & \textbf{2138012} \\ \hline
\end{tabularx}

% ==========================================
% 2. ORDINE DEL GIORNO
% ==========================================
\section{Ordine del Giorno}
Gli argomenti previsti per la riunione odierna sono:
\begin{enumerate}
    \item Determinare ruolo di responsabile
    \item Revisione valutazione capitolati
    \item Valutazioni way of working
    \item Revisione dichiarazione degli impegni
\end{enumerate}

% ==========================================
% 3. RESOCONTO DELLA RIUNIONE
% ==========================================
\section{Resoconto della Riunione}

\subsection{Determinata ruolo di responsabile}
Per il primo sprint è stato scelto come responsabile Riccardo Pozzato, al termine delle due settimane verrà sostituito da un'altro membro.

\subsection{Revisione valutazione capitolati}
Rivista la valutazione generale dei capitolati cambiando i vari voti assegnati per rispecchiare più accuratamente le preferenze del gruppo. Rivista la distribuzione delle ore lavoro per i membri, migliorata impaginazione e finalizzati pro, contro e sintesi.

\subsection{Way of Working}
È stata effettuata una riorganizzazione delle metodologie di lavoro del gruppo, con particolare attenzione alla gestione del controllo versione e delle procedure di integrazione:

\begin{itemize}
    \item \textbf{Adozione del Trunk-Based Development:} Transizione da GitFlow a Trunk-Based per superare la complessità degli allineamenti tra branch. La metodologia è stata successivamente adattata al contesto di progetto.
    \item \textbf{Protezione del ramo principale:} Il trunk è protetto da \textit{ruleset} che vietano modifiche dirette, richiedendo obbligatoriamente l'approvazione di 2 revisori e il superamento della pipeline di CI.
    \item \textbf{Tracciabilità e convenzioni:} Le Pull Request devono seguire regole di naming uniformi al messaggio dell'ultimo commit, conformemente alle convenzioni adottate per garantire un tracciamento chiaro di ogni modifica.
    \item \textbf{Gestione degli stack con SPR:} Utilizzo di Stacked PR per la gestione delle Pull Request. Per ovviare ai problemi di tracciamento remoto del tool, è stata implementata una soluzione interna basata su alias dedicati.
    \item \textbf{Criteri di accettazione:} Definizione e impiego di criteri di accettazione non binari per le attività.
\end{itemize}

\subsection{Dichiarazione degli impegni}
Finalizzata stesura della dichiarazione degli impegni includendo anche spiegazioni sul Way of working pianificato dal gruppo.



\subsection{Decisioni prese}
\begin{itemize}
    \item \textbf{V-1.1} Scelta responsabile primo sprint
    \item \textbf{V-1.2} Consolidato il Way of Working
    \item \textbf{V-1.3} Consolidata la dichiarazione degli impegni
\end{itemize}

% ==========================================
% 4. AZIONI DA SVOLGERE (TO-DO)
% ==========================================
\section{Azioni da Svolgere (To-Do)}
Nessun compito individuale da svolgere. Le attività elencate sono ricostruite dalle modifiche successive al repository.

\begin{tabularx}{\textwidth}{|c|X|l|l|}
\hline
\textbf{Codice Rif.} & \textbf{Descrizione Task} & \textbf{Assegnatario} & \textbf{Scadenza} \\ \hline
--- & Pubblicare la valutazione approfondita dei capitolati e la dichiarazione degli impegni, trattate durante la riunione. Entrambe sono state inserite nel repository il 30/10/2026 (\href{https://github.com/BitCrusher-Team/Progetto_SWE_26-27/pull/27}{PR \#27}). & --- & --- \\ \hline
--- & Pubblicare il verbale del 29/10/2026 con il riferimento nel README (\href{https://github.com/BitCrusher-Team/Progetto_SWE_26-27/pull/33}{PR \#33}, 30/10/2026) e aggiornare i percorsi della documentazione di candidatura sul sito (\href{https://github.com/BitCrusher-Team/Progetto_SWE_26-27/pull/31}{PR \#31}, 30/10/2026). & --- & --- \\ \hline

\end{tabularx}

% ==========================================
% 5. PROSSIMA RIUNIONE
% ==========================================
\section{Prossima Riunione}
\begin{itemize}
    \item \textbf{Data:} 2026-10-31
    \item \textbf{Ora:} 15:00
    \item \textbf{OdG preliminare:} PowerPoint per il prossimo Diario di Bordo
\end{itemize}

\end{document}